
\needspace{3\baselineskip}
\color{uniblue}{\subsection{Блокировки управление приводом РПН}}
\color{black}

\begin{enumerate}[label=\arabic{section}.\arabic{subsection}.\arabic{enumi}, labelsep=4pt, leftmargin=0pt, itemindent=57pt, itemsep=0pt, parsep=5pt] % если что можно крутить itemsep

%%%%%%%%%%%%%%%%%%%%%%%%%%%%%%%%%%%%%%%%%%%%%%%%%%%%%%%%%%%%%%%%%%%%% ОБЩИЕ СВЕДЕНИЯ %%%%%%%%%%%%%%%%%%%%%%%%%%%%%%%%%%%%%%%%%%%%%%%%%%%%%%%%%%%%%%

\item Общие сведения

\begin{enumerate}[label=\arabic{section}.\arabic{subsection}.\arabic{enumi}.\arabic*, labelsep=4pt, leftmargin=0em, itemindent=65pt, parsep=0pt]

\item В зависимости от режима работы сети, температуры, уровня масла в отсеке контактора РПН и состояния привода (отказ, местное управление, крайнее положение и др.) возможны ситуации, когда управление приводом должно быть блокированным. Для того, чтобы избежать возможного повреждения оборудования или возникновения нежелательного режима сети в устройстве предусмотрена функция блокировки управления. 
\item В зависимости от возникающей ситуации блокировка действует на команды управления в сторону увеличения напряжения, в сторону снижения или в обоих направлениях.

\item Запрет выдачи команды <<Убавить>> формируется в следующих случаях:
\begin{itemize}
    \item крайнее нижнее положение (от концевого выключателя привода);
    \item допустимое нижнее положение (блокирует только автоматические команды, см. \ref{tekush}).
\end{itemize}

\item Запрет выдачи команды <<Прибавить>> формируется в следующих случаях (блокирует только автоматические команды, см. \ref{tekush}):
\begin{itemize}
    \item перегрузка секции;
    \item высокое напряжение секции;
    \item низкое напряжение секции;
    \item повышение напряжения нулевой последовательности на секции;
    \item повышение напряжения обратной последовательности на секции;
    \item крайнее верхнее положение (от концевого выключателя привода, см. \ref{tekush});
    \item допустимое верхнее положение (блокирует только автоматические команды, см. \ref{tekush}).
\end{itemize}

\item Запрет выдачи обеих команд <<Убавить>> и <<Прибавить>> формируется в следующих случаях:
\begin{itemize}
    \item перегрузка контактора РПН (контроль по стороне ВН);
    \item снижение температуры масла;
    \item снижение уровня масла;
    \item привод в режиме местного управления;
    \item отказ привода.
\end{itemize}

\item
В таблице \ref{privod:tbl1} приведены параметры, необходимые для настройки блокировок управления приводом РПН.

\small
\begin{longtable}{|>{\centering\arraybackslash}m{5.3cm}|>{\centering\arraybackslash}m{3.3cm}|>{\centering\arraybackslash}m{4.2cm}|>{\centering\arraybackslash}m{1.8cm}|>{\centering\arraybackslash}m{1cm}|}
\caption{Параметры для настройки блокировок управления приводом РПН\hfill\vspace{-0.5\baselineskip}}\label{privod:tbl1}\\ 
\hline
\rowcolor{gray!30}
Параметр (Параметр на ИЧМ) & Условное обозначение на схеме & Значение/ Диапазон & Единица измерения & Шаг \\ 
\hline
\endfirsthead
\caption*{\hspace{3pt}\emph{Продолжение таблицы \ref{privod:tbl1}\hfill\vspace{-0.5\baselineskip}}} \\ % сделано по ГОСТ 2.105 п.6.8.7
\hline
\rowcolor{gray!30}
Параметр (Параметр на ИЧМ) & Условное обозначение на схеме & Значение/ Диапазон & Единица измерения & Шаг \\ 
%\hline
\endhead
%\hline
\endfoot
%\hline
\endlastfoot
%==+t1*AVTR1|_TTCC3
\centering Блокировка от контролируемой секции (БлкКонтрСекц) & \centering SGF1 & \centering 0 = Не предусмотрено\\1 = Предусмотрено & \centering -- & \centering \arraybackslash -- \\
\hline
%===t1
%==+t1*ATCC1|_TTCC3
\centering Допустимый ток перегрузки секции, выше которого запрет на повышение (Iвв>блк) & \centering Iвв>блк & \centering 1,000 ... 2,000 & \centering о.е. & \centering \arraybackslash 0,001 \\
\hline
\centering Верхняя уставка напряжения для блокировки автоматических команд на понижение (U\verb|>>|блк) & \centering U\verb|>>|блк & \centering 1,00 ... 2,00 & \centering о.е. & \centering \arraybackslash 0,01 \\
\hline
\centering Нижняя уставка напряжения блокировки автоматических команд на повышение (U\verb|<<|блк) & \centering U\verb|<<|блк & \centering 0,00 ... 1,00 & \centering о.е. & \centering \arraybackslash 0,01 \\
\hline
\centering Максимальное утроенное напряжения нулевой последовательности блокировки автоматических команд на повышение (3U0блк) & \centering 3U0блк & \centering 0,00 ... 2,00 & \centering о.е. & \centering \arraybackslash 0,01 \\
\hline
\centering Максимальное напряжение обратной последовательности блокировки автоматических команд на повышение (U2блк) & \centering U2блк & \centering 0,00 ... 2,00 & \centering о.е. & \centering \arraybackslash 0,01 \\
\hline
\centering Ток перегрузки РПН (Iрпн>блк) & \centering Iрпн>блк & \centering 1,00 ... 2,00 & \centering о.е. & \centering \arraybackslash 0,01 \\
\hline
%===t1
\end{longtable}
\normalsize


\end{enumerate}

%%%%%%%%%%%%%%%%%%%%%%%%%%%%%%%%%%%%%%%%%%%%%%%%%%%%%%%%%%%%%%%%%%%%% Контроль секции %%%%%%%%%%%%%%%%%%%%%%%%%%%%%%%%%%%%%%%%%%%%%%%%%%%%%%%%%%%%%%

\item Контроль режима работы секций

\begin{enumerate}[label=\arabic{section}.\arabic{subsection}.\arabic{enumi}.\arabic*, labelsep=4pt, leftmargin=0em, itemindent=65pt, parsep=0pt]

\item
Устройство контролирует значение напряжения, уровень несимметрии напряжения, нагрузку на регулируемой и контролируемой (при наличии) секции для блокировки команд управления в автоматическом режиме. Для каждой секции предусмотрены ИО для контроля повышения, снижения напряжения, повышения напряжения обратной и нулевой последовательности, перегрузки секции (см. рис. \ref{privod:img1}).  
\item
Возможность ввода блокировки автоматических команд управления от ИО контролируемой секции определяется положением программного переключателя SGF1. На рисунке \ref{privod:img2} приведена схема логики блокировки управления приводом по режиму работы сети. 

\vspace{2mm}
\begin{figure}[H]
\centering
\includegraphics[width=0.67\textwidth,height=0.67\textheight,keepaspectratio]{\fbpath/503. Привод РПН/_latex/img/1.pdf}
\captionof{figure}{Функциональная схема измерительных органов логики формирования сигналов нарушения нормальной работы секций}\label{privod:img1}
\end{figure}

\vspace{2mm}
\begin{figure}[H]
\centering
\includegraphics[width=1\textwidth,height=1\textheight,keepaspectratio]{\fbpath/503. Привод РПН/_latex/img/2.pdf}
\captionof{figure}{Функциональная схема логики блокировки управления приводом по режиму работы сети}\label{privod:img2}
\end{figure}

\end{enumerate}

%%%%%%%%%%%%%%%%%%%%%%%%%%%%%%%%%%%%%%%%%%%%%%%%%%%%%%%%%%%%%%%%%%%%% Контроль перегрузки по стороне ВН трансформатора %%%%%%%%%%%%%%%%%%%%%%%%%%%%%%%%%%%%%%%%%%%%%%%%%%%%%%%%%%%%%%

\item Контроль перегрузки по стороне ВН трансформатора

\begin{enumerate}[label=\arabic{section}.\arabic{subsection}.\arabic{enumi}.\arabic*, labelsep=4pt, leftmargin=0em, itemindent=65pt, parsep=0pt]

\item    
Чтобы исключить возможность повреждения контактора РПН при коммутациях  высоких значений тока, в устройстве предусмотрена функция запрета регулирования при перегрузке по стороне ВН трансформатора (см. рис. \ref{privod:img3}). 
\item 
Сигнал о перегрузке  может поступать в устройство через дискретный вход (для исполнения I) или GOOSE сообщение (для исполнений I и II), или формироваться в устройстве при срабатывании специального ИО тока стороны ВН трансформатора (для исполнения II). При появлении сигнала о перегрузке немедленно формируется сигнал запрета управления приводом во всех режимах управления.

\vspace{2mm}
\begin{figure}[H]
\centering
\includegraphics[width=0.7\textwidth,height=0.7\textheight,keepaspectratio]{\fbpath/503. Привод РПН/_latex/img/3.pdf}
\captionof{figure}{Функциональная схема логики контроля тока стороны ВН трансформатора}\label{privod:img3}
\end{figure}

\end{enumerate}

%%%%%%%%%%%%%%%%%%%%%%%%%%%%%%%%%%%%%%%%%%%%%%%%%%%%%%%%%%%%%%%%%%%%% Запрет управления по температуре масла %%%%%%%%%%%%%%%%%%%%%%%%%%%%%%%%%%%%%%%%%%%%%%%%%%%%%%%%%%%%%%

\item Запрет управления по температуре масла
\begin{enumerate}[label=\arabic{section}.\arabic{subsection}.\arabic{enumi}.\arabic*, labelsep=4pt, leftmargin=0em, itemindent=65pt, parsep=0pt]
\item 
При низкой ($t<-20\ldots25~^\circ$C) температуре окружающей среды многократно увеличивается вязкость масла, что затрудняет процесс переключения контактора. При большой вязкости масла может произойти повреждение токоограничивающих элементов в связи с увеличением времени переключения или поломка контактора в связи с увеличением механической нагрузки.
\item
Сигнал о снижении температуры масла поступает в устройство от внешнего датчика посредством GOOSE сообщения (исполнение I и II) или дискретный вход (исполнение I). Получив данный сигнал, логика устройства блокирует команды управления приводом. 
\end{enumerate}

%%%%%%%%%%%%%%%%%%%%%%%%%%%%%%%%%%%%%%%%%%%% Запрет управления по конечным выключателям или допустимым крайним положениям %%%%%%%%%%%%%%%%%%%%%%%%%%%%%%%%%%%%%%%%%%%%%%%%%%

\item Запрет управления по концевым выключателям или допустимым крайним положениям
\begin{enumerate}[label=\arabic{section}.\arabic{subsection}.\arabic{enumi}.\arabic*, labelsep=4pt, leftmargin=0em, itemindent=65pt, parsep=0pt]
\item 
Устройство отслеживает положения концевых выключателей привода. При достижении крайнего верхнего или крайнего нижнего положений формируются соответствующие сигналы <<Запрет прибавить>> или <<Запрет убавить>>.
\item
Кроме того в устройстве предусмотрена возможность ограничения уставками диапазона регулирования в автоматическом режиме. При достижении крайнего верхнего или нижнего допустимого положения при автоматическом регулировании логикой устройства формируются соответствующие сигналы <<Запрет прибавить>> или <<Запрет убавить>>.
\end{enumerate}

%%%%%%%%%%%%%%%%%%%%%%%%%%%%%%%%%%%%%%%%%%%% Запрет управления по сигналам неисправности привода и трансформатора %%%%%%%%%%%%%%%%%%%%%%%%%%%%%%%%%%%%%%%%%%%%%%%%%%

\item Запрет управления по сигналам неисправности привода и трансформатора
\begin{enumerate}[label=\arabic{section}.\arabic{subsection}.\arabic{enumi}.\arabic*, labelsep=4pt, leftmargin=0em, itemindent=65pt, parsep=0pt]
\item Кроме указанных выше сигналов о снижении температуры масла и о достижении приводом крайних положений устройство может контролировать следующие сигналы с привода РПН и с силового трансформатора:
\begin{itemize}
    \item отключено автоматический выключатель питание привода;
    \item минимальный уровень масла в отсеке контактора РПН;
    \item привод в режиме местного управления.
\end{itemize}
\item 
Сигналы могут поступать в устройство посредством GOOSE сообщения (исполнение I и II) или дискретный вход (исполнение I). Логикой устройства в процессе переключения может формироваться сигнал «ПМ отказал».
\item
Указанные выше  сигналы блокируют команды управления приводом во всех режимах управления. Схема выходных цепей функции блокировки управления приводом приведена на рисунке \ref{privod:img4}.
\end{enumerate}

\vspace{2mm}
\begin{figure}[H]
\centering
\includegraphics[width=0.6\textwidth,height=0.6\textheight,keepaspectratio]{\fbpath/503. Привод РПН/_latex/img/4.pdf}
\captionof{figure}{Выходные цепи функции блокировки управления приводом}\label{privod:img4}
\end{figure}

\end{enumerate}